\documentclass[a4paper, 11pt]{article}

\usepackage[T1,T2A]{fontenc}
\usepackage[utf8]{inputenc}
\usepackage[english, russian]{babel}

\usepackage{comment}
\usepackage{fullpage}
\usepackage{amssymb}
\usepackage{amsmath}
\usepackage{amsthm}
\usepackage{enumitem}

% %\DeclareMathOperator{\ord}{ord}
\DeclareMathOperator{\chr}{char}
\DeclareMathOperator{\Mod}{mod}
\DeclareMathOperator{\Gal}{Gal}
\DeclareMathOperator{\N}{N}
\DeclareMathOperator{\Tr}{Tr}
\DeclareMathOperator{\SL}{SL}

\begin{document}

\begin{center} \Large \bf Листок 6\\ Тема 6 (2.2). Расширения конечных полей \end{center}

\noindent\textbf{Упражнения и задачи}

\begin{enumerate}[topsep=0pt]
    
    \item Пусть $n|q-1$, $F=\mathbb{F}_q$, $K/F$ --- расширение конечных полей, $[K:F] = n$. Докажите, что $\forall \alpha \in F^*$ уравнение $x^n=\alpha$ имеет $n$ решений в $K$. %[IR, p86 ex7.5]
    \item Пусть $K/F$ --- расширение конечных полей, $\mathrm{char} (F) \neq 2$, $[K:F] = 3$, $\alpha \in F$. Докажите, что если $\alpha$ не является квадратом в $F$, то $\alpha$ не является квадратом и в $K$. %[IR, p86 ex7.6]
    \item Пусть $F=\mathbb{F}_q$, $K/F$ --- расширение конечных полей, $\alpha \in \mathbb{F}_q$, $n|q-1$ и $x^n=\alpha$ не разрешимо в $\mathbb{F}_q$. Докажите, что $x^n=\alpha$ не разрешимо в $K$, если $(n,[K:F])=1$. %[IR, p86 ex7.9]
    \item Пусть $F=\mathbb{F}_q$, $K/F$ --- расширение конечных полей, $[K:F] = 2$. Докажите, что $\forall \beta \in K$ $\beta^{1+q} \in F$. Более того, $\forall \alpha \in F$ $\exists \beta \in K$: $\alpha=\beta^{1+q}$. %[IR, p86 ex7.10]
    \item Докажите, что число упорядоченных базисов $\mathbb{F}_{q^m}$ над $\mathbb{F}_q$ равно $(q^m-1)(q^m-q)\dots (q^m-q^{m-1})$. %[LN Ex 2.37]
    \item Пусть $\mathbb{F}_{q^m}/\mathbb{F}_q$ --- расширение конечных полей, $\alpha \in \mathbb{F}_{q^m}^*$. Докажите, что сопряженные элементы $\alpha, \alpha^q, \alpha^{q^2}, \dots, \alpha^{q^{m-1}}$ имеют один и тот же порядок в мультипликативной группе $\mathbb{F}_{q^m}^*$.  %[LN, p50 th2.18]
    \item Пусть $\mathbb{F}_{q^m}/\mathbb{F}_q$ --- расширение. Докажите, что $\forall c\in \mathbb{F}_q$
    $$
        \sum_{j=0}^{m-1} x^{q^j} - c = \prod_{\alpha\in \mathbb{F}_{q^m},\ \mathrm{Tr}_{\mathbb{F}_{q^m}/\mathbb{F}_q}(\alpha)=c} (x-\alpha).
    $$ %[LN Ex 2.34]
    \item Пусть $\mathbb{F}_{q^m}/\mathbb{F}_q$ --- расширение, $L$ --- $\mathbb{F}_q$-линейный оператор на $\mathbb{F}_{q^m}$ (как на векторном пространстве над $\mathbb{F}_q$). Докажите, что $\exists !\, \alpha_0, \dots, \alpha_{m-1}\in \mathbb{F}_{q^m}$: оператор $L$ имеет вид $L(\beta) = \alpha_0 \beta + \alpha_1 \beta^q + \dots + \alpha_{m-1}\beta^{q^{m-1}}$. %[LN Ex 2.36]
    \item Пусть $\mathbb{F}_{q^m}/\mathbb{F}_q$ --- расширение.
    \begin{itemize}[noitemsep,topsep=0pt]
        \item Докажите, что $\mathrm{Tr}_{\mathbb{F}_{q^m}/\mathbb{F}_q}(\alpha)=0$ $\Leftrightarrow$ $\alpha=\beta^q-\beta$ для некоторого $\beta\in\mathbb{F}_{q^m}$.
        \item ($^*$) Пусть $q=p^s$, $a\in\mathbb{F}_q$. Докажите, что $x^p-x-a$ неприводим над $\mathbb{F}_q$ $\Leftrightarrow$ $\mathrm{Tr}_{\mathbb{F}_q/\mathbb{F}_p}(a)\neq 0$.
    \end{itemize} %[LN Th 2.25, Th 3.78]
\end{enumerate}

% \noindent\textbf{SageMath}
% \begin{itemize}[topsep=0pt]
%     \item Исследуйте основные функции SageMath связанные с автоморфизмами и расширениями конечных полей:
%     \begin{itemize}[noitemsep,topsep=0pt]
%         \item Группа автоморфизмов поля: \texttt{End()};
%         \item Автоморфизм Фробениуса: \texttt{frobenius\_endomorphism()};
%         \item Базисы расширений конечных полей;
%         \item Характеристический многочлен \texttt{charpoly()};
%         \item Норма, след, дискриминант.
%      \end{itemize}
% \end{itemize}

\noindent\textbf{Темы для самостоятельного изучения}
\begin{itemize}[topsep=0pt]
    \item Нормальный базис и теорема о нормальном базисе, [LN].
\end{itemize}    

\end{document}
